% Notes of the Montreal seminar, 1970 (cote 129, typed part): archive cote 129, site page 56 (PDF page 57).
% typescript. Draft compared: none.
% AI-transcribed by Claude Opus 5.5 (claude-opus-5-5, adaptive thinking, effort high), Message Batches API, batch mtl129_full (msgbatch_01AGnDKHwy6ATCM7ehUa6263), 2026-11-13; not yet checked by eye.
\scanpage{129/56}
\mspage{\hand{III-1}}

\uline{GROUPES DE BARSOTTI-TATE}

1. \uline{Notations et définitions préléminaires}

Soit $S$ un schéma. On munit la catégorie Sch/S des schémas
au dessus de $S$ d'une topologie \hand{$\mathcal{T}$} que l'on suppose plus fine que la
topologie f.p.p.f. (fidèlement plate et de présentation finie) et moins
fine que la topologie f.p.q.c. (fidèlement plate et quasi-compacte)
(SGA 3 IV, 6.3). L'introduction de cette topologie est technique et n'est
pas essentielle dans les énoncés finaux (voir 5.1).

Pour abréger, nous dirons ici que $G$ est un groupe si $G$ est
un faisceau en groupes commutatifs sur Sch/S pour la topologie \hand{$\mathcal{T}$}.

Pour tout entier $n \geq 1$ et pour tout groupe $G$, on pose:
\[ \Lambda_n = \mathbb{Z}/p^n\mathbb{Z} \]
\[ G(n) = \mathrm{Ker}\,\{G \xrightarrow{p^n \mathrm{Id}} G\} \ , \]
et
\[ G(\infty) = \varinjlim_{n} G(n) \ ; \]
$G(n)$ est un $\Lambda_n$-module (i.e. un \hand{$\mathcal{T}$}-faisceau en $\Lambda_n$-modules);
$G(\infty)$ est le plus grand sous-groupe de $p$-torsion de $G$. On a évidemment:
\[ G(n)(n') = G(n') \qquad \text{pour} \quad n' \leq n \leq \infty \]

Si l'on veut indiquer qu'un groupe est annulé par $p^n$, c'est-à-%
dire est un $\Lambda_n$-module, on pourra donc le noter $G(n)$ au lieu de $G$ sans
entraîner de confusions.

\mspage{\hand{56}}
